%قبل از شروع مسئله ترجیحاً چیزی ننویسید. اما اگر خواستید بنویسید، در قالب لیستی (فقرات و شمارش) بنویسید.
%برای اضافه کردن نمره، بعدی از شروع‌مسئله، [x] را اضافه کنید و به جای x  نمره مورد نظر را قرار دهید. برای زیرمسئله‌ها هم به همین صورت است.
%\شروع‌مسئله[10]
\شروع‌مسئله [12]
%\متن‌سیاه{سوال عملی اول:}\\
%\متن‌سیاه{دراپ‌تست}\\
محمد معلم کلاسی متشکل از \متن‌لاتین{n} دانش‌آموز است. محمد بچه های کلاس را کنار هم در یک صف قرار می‌دهد. او که شدیدا به نظم اهمیت می‌دهد، می‌خواهد بچه ها را بر حسب قدشان از چپ به راست،‌ به شکل نزولی مرتب کند. برای این کار هر بار می‌تواند دو دانش‌آموز مجاور را با هم جا به جا کند. اما او که خیلی کار دارد و حسابی خسته است،‌ می‌خواهد با کمترین تعداد حرکات این کار را انجام دهد. برای این کار هم از شما خواسته است که مینیمم تعداد حرکات را بدست آورید.\\
دقت کنید که در این سوال از کد های آماده سایت‌ها و توابع مرتب سازی آماده نمی‌توانید استفاده کنید. هم چنین باید از مباحث آموخته شده در درس استفاده کنید.

\متن‌سیاه{ورودی:}\\
در خط اول عدد \متن‌لاتین{n} می‌آید که تعداد دانش‌آموز هاست. در خط بعدی \متن‌لاتین{n} عدد \متن‌لاتین{a[i]} پشت سر هم می‌آید که نشان دهنده‌ی قد افراد کلاس است.\\
\متن‌سیاه{خروجی:}\\
یک عدد خروجی دهید که نشان‌دهنده‌ی کمترین تعداد حرکات لازم است.\\
\متن‌سیاه{محدودیت‌ها:}\\
$1 \leq n \leq 10^{5}$\\
$0 \leq a[i] \leq 10^{9}$\\

\متن‌سیاه{ورودی و خروجی نمونه:}\\

تست اول
\begin{latin}
	\begin{tabularx}{\textwidth}{|X | X|}
		\hline
		\textbf{stdin} &
		\textbf{stdout}\\
		\hline
		4 & 4\\
		0 3 4 2&\\
		\hline
	\end{tabularx}
\end{latin}

تست دوم
\begin{latin}
	\begin{tabularx}{\textwidth}{|X | X|}
		\hline
		\textbf{stdin} &
		\textbf{stdout}\\
		\hline
		5 & 1\\
		5 2 4 2 1&\\
		\hline
		
	\end{tabularx}
\end{latin}
\solution{
	می‌خواهیم قدهای افراد را به صورت نزولی مرتب کنیم. با جابه‌جا کردن دو نفر، تعداد نابه‌جایی هایی که قد افراد می‌سازد حداکثر یکی تغییر می‌کند. پس می‌دانیم برای رسیدن به حالت نهایی، لازم است حداقل به اندازه‌ی اختلاف نابه‌جایی های حالت اولیه و حالت ثانویه هزینه کنیم. هم‌چنین راهی با این تعداد حرکت وجود دارد. زیرا تا وقتی به حالت نهایی نرسیده‌ایم، دو نفر کنار هم هستند که قدشان نا‌به‌جایی نمی‌سازد و می‌توانیم این دو نفر را جابه‌جا کنیم و یکی به تعداد نابه‌جایی ها اضافه شود(درستی این فرض را می‌توانید به سادگی نشان دهید). برای بدست آوردن تعداد نابه‌جایی ها نیز می‌توانیم از مرتب‌سازی ادغامی استفاده کنیم. بدین شکل که به صورت بازگشتی تعداد نابه‌جایی ها در نیمه‌ی چپ و راست را بدست آوریم و در ادغام این دو قسمت، هنگامی که یک عدد از یک نیمه انتخاب می‌شود ، با توجه به تعداد اعدادی قبل‌تر که از نیمه‌ی دیگر انتخاب شده است،‌ می‌فهمیم چند نابه‌جایی با اعداد نیمه‌ی دیگر می‌سازد. اردر این راه نیز اردر مرتب‌سازی ادغامی یعنی
$O(nlogn)$
 خواهد بود. در زیر کد این سوال را که توسط آقای علی جوانمرد زده شده است می‌بینید:

}
%}
% صورت اصلی مسئله را در این‌جا بنویسید.
%اگر مسئله چند قسمت داشت به این صورت عمل کنید:
%\شروع{زیرمسائل}
%\زیرمسئله[0] % مسئلۀ ۱آ
%\زیرمسئله % مسئلۀ ۱ب
%\زیرمسئله % مسئلۀ ۱پ
%\زیرمسئله % مسئلۀ ۱ت
%\پایان{زیرمسائل}



%%%%%%%%%%%%%%%%%%%%%%%%%%%%%%%%%%%%%%%%%%%%%%%%%%%%%
% برخی دستورات مفید برای تمرین‌ها. %
%%%%%%%%%%%%%%%%%%%%%%%%%%%%%%%%%%%%%%%%%%%%%%%%%%%%%

% فرمول‌های نمونه:
%$f(x) = \Theta(x)$
%$T(x, y) \leq \log(x) + 2^y + \binom{2n}{n}$
% {\tt code\_function}


% لیست موارد بدون شماره‌گذاری:
%\شروع{فقرات}
%    \فقره مورد اول 
%        \شروع{فقرات}
%            \فقره مورد اول داخلی 
%                \شروع{فقرات}
%                    \فقره اولین مورد خیلی داخلی 
%                    \فقره دومین مورد خیلی داخلی 
%                \پایان{فقرات}
%            \فقره مورد دوم داخلی 
%        \پایان{فقرات}
%    \فقره مورد دوم 
%\پایان{فقرات}

% لیست موارد شماره‌گذاری شده:
%\شروع{شمارش}
%    \فقره یک مورد 
%    \فقره یک مورد دیگر 
%    \فقره باز هم یک مورد دیگر
%\پایان{شمارش}

% برای نوشتن یک فرمول چند خطی که یک شماره بخوره. برای حالت‌های دیگر به سایت زیر مراجعه کنید:
% https://www.overleaf.com/learn/latex/Aligning_equations_with_amsmath
%\begin{align} 
%    \begin{split}
%        (x+y)^3 &= (x+y)^2(x+y) \\
%                &= (x^2+2xy+y^2)(x+y) \\
%                &= (x^3+2x^2y+xy^2) + (x^2y+2xy^2+y^3) \\
%                &= x^3+3x^2y+3xy^2+y^3
%    \end{split}                                 
%\end{align}

% نوشتن ماتریس و آرایه دوبعدی
%\begin{align}
%    A = 
%    \begin{bmatrix}
%        A_{11} & A_{21} \\
%        A_{21} & A_{22}
%    \end{bmatrix}
%\end{align}

% اضافه کردن تصویر با شرح. برای اضافه کردن تصویر بدون شرح، فقط دستور درج‌تصویر را استفاده کنید.
%\شروع{شکل}[h!]
%\تنظیم‌ازوسط
%  \درج‌تصویر[پهنا=0.5\پهنای‌متن]{img.jpg}
%\شرح{شرح تصویر}
%\برچسب{تصویر:۱}
%\پایان{شکل}
%\پایان{زیرمسائل}

%کد زیر را ببینید.
%%\lstinputlisting{source_filename.py}
\begin{latin}
\begin{lstlisting}
def merge_sort_slice(array):
    if (len(array) <= 1):
        return array
    m = len(array) // 2
    L = merge_sort_slice(array[:m])
    R = merge_sort_slice(array[m:])
    return list(merge(L, R))


def merge(L, R):
    global b
    l = r = 0
    while l < len(L) or r < len(R):
        if l < len(L) and (
                r >= len(R) or L[l] >= R[r]):
            yield L[l]
            l += 1
        else:
            b += len(L) - l
            yield R[r]
            r += 1


number = int(input())
b = 0
s = []
s = list(map(int, input().split()))
s = merge_sort_slice(s)
print(b)
\end{lstlisting}
\end{latin}

